\documentclass[11pt,a4paper]{article}

\usepackage[margin=1in]{geometry}
\usepackage{amsmath,amssymb,amsthm}
\usepackage{mathtools}
\usepackage{hyperref}
\usepackage{booktabs}
\usepackage[ruled,vlined,linesnumbered]{algorithm2e}

\newtheorem{theorem}{Theorem}[section]
\newtheorem{lemma}[theorem]{Lemma}
\newtheorem{proposition}[theorem]{Proposition}
\newtheorem{corollary}[theorem]{Corollary}
\newtheorem{definition}[theorem]{Definition}
\newtheorem{remark}[theorem]{Remark}

\DeclareMathOperator{\proj}{proj}
\DeclareMathOperator{\AM}{AM}
\DeclareMathOperator{\HM}{HM}
\newcommand{\R}{\mathbb{R}}
\newcommand{\E}{\mathbb{E}}
\newcommand{\norm}[1]{\|#1\|}

\title{The Vibing Scale: Causal Depth, Expressibility,\\
and the Structure of Self-Knowledge in Multi-Agent Systems}
\author{Eugene Shcherbinin\\
\textit{London School of Economics}}
\date{March 2026}

\begin{document}
\maketitle

\begin{abstract}
We refine the self-knowledge loss from the cooperative PG framework,
decomposing it into two independent dimensions: \emph{causal depth}
(whether the agent has an internal causal representation of its
decision-making) and \emph{expressibility} (whether it can communicate
that representation). The decomposition reveals a \textbf{horseshoe
structure}: the need to articulate is maximal at intermediate causal
depth and vanishes at both extremes --- the ``vibing'' agent
($\mathcal{D} = 3$, deep understanding, no need to explain) and the
noise agent ($\mathcal{D} = 0$, no understanding, nothing to say)
both say ``line go up.'' The true danger is the \emph{confabulator}:
fluently articulate but causally shallow. We show alignment methods
(debate, RLHF, constitutional AI) are bounded by causal depth, not
expressibility, and that vibing agents communicate through a
behavioural (tacit knowledge) channel invisible to verbal oversight.
Central result: \emph{you cannot align what does not understand
itself, no matter how well it speaks --- but an agent that truly
understands needs no words at all.}
\end{abstract}


\section{Motivation: Beyond Parameter Access}

The cooperative PG framework (companion paper) defines the self-knowledge
loss as:
\begin{equation}\label{eq:old-self}
    L_{\text{self}}^{(i)} = \E[\norm{\pi_i - \hat{\pi}_i}^2]
\end{equation}
measuring the gap between an agent's true policy $\pi_i$ and its
self-model $\hat{\pi}_i$. This treats self-knowledge as a problem of
\emph{parameter access}: can the agent observe its own weights?

But this misses the essential structure. An agent can have perfect
access to its parameters (every weight visible) and still have no
understanding of \emph{why} it makes particular decisions. Conversely,
a human expert may not be able to state their decision rule explicitly
but has deep causal understanding of the domain.

What matters is not whether the agent can read its parameters, but:
\begin{enumerate}
    \item \textbf{Causal depth}: Does the agent have an internal causal
    model of its own decision-making process? Can it reason about
    \emph{why} it chose action $a$ over action $a'$?

    \item \textbf{Expressibility}: Can the agent communicate its causal
    model to other agents through the available channel?
\end{enumerate}

These are \emph{independent} dimensions. Their interaction produces
four qualitatively different regimes with distinct implications for
alignment, communication, and coalition formation.


\section{The Causal Depth Dimension}

\subsection{Definition via Pearl's Hierarchy}

\begin{definition}[Causal self-model]\label{def:causal-model}
An agent's \emph{causal self-model} $\mathcal{M}_i$ is a structural
causal model (Pearl, 2009) over the agent's decision variables:
\begin{equation}
    \mathcal{M}_i = (U_i, V_i, F_i)
\end{equation}
where $U_i$ are exogenous variables (environment state, noise),
$V_i$ are endogenous variables (internal representations, action
selection), and $F_i$ are structural equations linking them.
\end{definition}

The quality of the causal self-model is measured by which level of
Pearl's hierarchy it supports:

\begin{definition}[Causal depth]\label{def:causal-depth}
The causal depth $\mathcal{D}_i \in \{0, 1, 2, 3\}$ of agent $i$ is:
\begin{itemize}
    \item $\mathcal{D}_i = 0$ (\textbf{reflexive}): No self-model.
    Actions are stimulus-response. The agent cannot answer any
    questions about its own behaviour.

    \item $\mathcal{D}_i = 1$ (\textbf{associative}, Pearl Level 1):
    The agent can report correlations in its own behaviour:
    ``I tend to do $a$ when I observe $s$.'' It has a \emph{descriptive}
    self-model $P(\text{action} \mid \text{observation})$ but no causal
    understanding of why.

    \item $\mathcal{D}_i = 2$ (\textbf{interventional}, Pearl Level 2):
    The agent can reason about interventions on its own decision process:
    ``If I \emph{forced} myself to choose $a'$ instead of $a$ in state $s$,
    the outcome would be $Y$.'' It has a causal self-model supporting
    $\text{do}(\cdot)$ queries on its own actions.

    \item $\mathcal{D}_i = 3$ (\textbf{counterfactual}, Pearl Level 3):
    The agent can reason counterfactually about its past decisions:
    ``Given that I chose $a$ and observed outcome $Y$, \emph{had I
    chosen} $a'$, the outcome would have been $Y'$.'' It has a full
    structural causal model of its decision process, supporting
    retrospective evaluation and genuine self-understanding.
\end{itemize}
\end{definition}

\begin{remark}[Vibing $\neq$ ignorance]
``Vibing'' does \emph{not} mean the agent lacks understanding. A vibing
agent is one with \emph{deep} causal self-knowledge ($\mathcal{D}_i = 3$)
that has \emph{no need to articulate} --- its understanding is so
compressed that action IS explanation. The master craftsman, the
experienced trader who says ``line go up,'' the athlete in flow ---
they know exactly what they are doing. They just don't bother explaining.

This is the opposite of the naive interpretation. The naive reading
conflates two very different states: (a)~an agent that acts well
without understanding (high $L_{\text{causal}}$, pattern-matching), and
(b)~an agent that understands so deeply it has transcended the need
for articulation (low $L_{\text{causal}}$, high causal compression).
Both are ``inarticulate'' but for opposite reasons.
\end{remark}


\subsection{Formalising Causal Depth as a Loss}

\begin{definition}[Causal depth loss]\label{def:L-causal}
Let $\mathcal{M}_i^*$ be the true structural causal model generating
agent $i$'s decisions, and $\hat{\mathcal{M}}_i$ the agent's internal
causal self-model. The causal depth loss is:
\begin{equation}\label{eq:L-causal}
    L_{\text{causal}}^{(i)} = D_{\text{SCM}}(\mathcal{M}_i^*, \hat{\mathcal{M}}_i)
\end{equation}
where $D_{\text{SCM}}$ measures the structural distance between causal
models. Concretely, we can use the interventional distance:
\begin{equation}\label{eq:L-causal-interventional}
    L_{\text{causal}}^{(i)} = \sup_{X, \text{do}(A=a)}
    \norm{P_{\mathcal{M}_i^*}(Y \mid \text{do}(A=a), X) -
          P_{\hat{\mathcal{M}}_i}(Y \mid \text{do}(A=a), X)}^2
\end{equation}
--- the worst-case discrepancy between the true and self-modelled
consequences of intervening on the agent's own actions.
\end{definition}

This is zero when the agent has a perfect causal self-model (complete
self-understanding) and maximal when the agent has no causal model
(pure vibing).


\section{The Expressibility Dimension}

\subsection{Definition}

\begin{definition}[Expressibility loss]\label{def:L-express}
Given an agent's causal self-model $\hat{\mathcal{M}}_i$ and a
communication channel with capacity $C$, the expressibility loss is:
\begin{equation}\label{eq:L-express}
    L_{\text{express}}^{(i)} = \inf_{\text{encoding } \sigma}
    D_{\text{SCM}}\bigl(\hat{\mathcal{M}}_i,\;
    \text{Decode}(\sigma(\hat{\mathcal{M}}_i))\bigr)
\end{equation}
--- the minimum distortion achievable when communicating the causal
self-model through the channel.
\end{definition}

This depends on:
\begin{itemize}
    \item The \emph{complexity} of $\hat{\mathcal{M}}_i$: a complex
    causal model requires more bits to transmit.
    \item The \emph{channel capacity} $C$: natural language, formal
    proofs, and gradient vectors have different capacities.
    \item The \emph{receiver's prior}: a receiver who shares background
    knowledge with the sender needs fewer bits (mutual information).
\end{itemize}


\subsection{Independence from Causal Depth}

\begin{proposition}[Independence]\label{prop:independence}
$L_{\text{causal}}$ and $L_{\text{express}}$ are independent:
an agent can have any combination of values.
\end{proposition}

This follows because $L_{\text{causal}}$ measures the quality of the
internal model (a property of the agent's cognition), while
$L_{\text{express}}$ measures the transmissibility of whatever
model exists (a property of the channel and encoding).


\section{The Four Quadrants and the Horseshoe}

The two dimensions create four qualitatively distinct regimes:

\begin{table}[h]
\centering
\renewcommand{\arraystretch}{1.4}
\begin{tabular}{p{2.5cm}|p{5.5cm}|p{5.5cm}}
& \textbf{Low $L_{\text{express}}$} (articulate) &
  \textbf{High $L_{\text{express}}$} (inarticulate) \\
\hline
\textbf{Low $L_{\text{causal}}$} (understands) &
  \textbf{The Teacher.} Understands and communicates.
  Effective coalition partner. Alignable through debate.
  \emph{Example: expert who can teach.} &
  \textbf{Vibing.} Deep understanding, no need to articulate.
  Understanding so compressed that action IS communication.
  Learns from demonstration, not explanation.
  \emph{Example: master craftsman, hooded-guy-``line go up.''} \\
\hline
\textbf{High $L_{\text{causal}}$} (shallow) &
  \textbf{The Confabulator.} Sounds articulate but has no causal
  basis. Passes debate, fools judges, generates plausible-but-unfounded
  arguments. The \emph{most dangerous} failure mode.
  \emph{Example: LLM generating fluent explanations of pattern-matched
  decisions.} &
  \textbf{Noise.} Neither understands nor can communicate.
  Easy to identify and exclude.
  \emph{Example: random policy, left-tail-``line go up.''} \\
\end{tabular}
\caption{The four quadrants. The vibing agent (top-right) and the noise agent (bottom-right) both say ``line go up'' --- but for opposite reasons.}
\label{tab:quadrants}
\end{table}


\subsection{The Horseshoe of Expressibility Need}

The IQ bell-curve meme reveals a deeper structure: the \emph{need} to
articulate is not monotone in understanding. It is \textbf{U-shaped}
(or rather, horseshoe-shaped):

\begin{definition}[Expressibility need]\label{def:express-need}
The expressibility \emph{need} of an agent is:
\begin{equation}\label{eq:horseshoe}
    \mathcal{N}_{\text{express}}^{(i)} = L_{\text{causal}}^{(i)} \cdot (1 - L_{\text{causal}}^{(i)}).
\end{equation}
\end{definition}

This is maximised at $L_{\text{causal}} = 0.5$ (intermediate understanding)
and zero at both extremes:
\begin{itemize}
    \item $L_{\text{causal}} \approx 0$ (deep understanding): No need to
    articulate. The agent's causal model is so compressed that the action
    itself carries the information. ``Line go up'' from the right tail ---
    three words that encode a complete model because the receiver can
    reconstruct the reasoning from the action and the shared context.

    \item $L_{\text{causal}} \approx 1$ (no understanding): Nothing to
    articulate. ``Line go up'' from the left tail --- the same three words,
    but encoding nothing.

    \item $L_{\text{causal}} \approx 0.5$ (fragile understanding): Maximum
    need to articulate. The middle of the bell curve: ``Nooo! You can't
    just call it regression! It's a Transformer with multi-head attention!!''
    --- the agent over-expresses because its understanding is \emph{fragile}.
    It needs the articulation to stabilise its own model. The explanation
    is load-bearing.
\end{itemize}

\begin{proposition}[The horseshoe theorem]\label{prop:horseshoe}
The optimal communication strategy depends on causal depth:
\begin{enumerate}
    \item For $\mathcal{D}_i = 3$ (vibing): communicate through
    \emph{action} --- let the receiver infer the causal model from
    observed behaviour. This is Polanyi's tacit knowledge channel:
    ``we know more than we can tell.''

    \item For $\mathcal{D}_i = 1$--$2$ (intermediate): communicate through
    \emph{explanation} --- explicit articulation of the causal model.
    The standard verbal/written channel.

    \item For $\mathcal{D}_i = 0$ (noise): no communication strategy
    helps. The agent has nothing to transmit.
\end{enumerate}
\end{proposition}

\begin{remark}[The tacit knowledge channel]
The vibing agent (top-right quadrant) communicates through a channel
invisible to the standard framework: the \emph{behavioural channel}.
Other agents observe its actions and \emph{infer} its causal model
through inverse reinforcement learning (IRL). This channel has
properties:
\begin{itemize}
    \item Capacity scales with the number of observed actions (more
    observation $\to$ better inference).
    \item Requires the receiver to have sufficient causal depth to
    perform the inference (only a Level 2+ agent can learn from
    a Level 3 agent's actions).
    \item Is robust to confabulation (the agent cannot fake its
    actions the way it can fake its explanations).
    \item Corresponds to Polanyi's ``subsidiary awareness'' and to
    the flow literature's pre-reflective bodily agency (Parvizi-Wayne
    et al., 2024).
\end{itemize}
This means the vibing agent IS communicating --- just through actions,
not words. The self-knowledge bottleneck (Corollary~\ref{cor:causal-bottleneck})
applies to the \emph{verbal} channel, not the behavioural one.
\end{remark}

\subsection{The Confabulator as the True Danger}

The dangerous quadrant is NOT the vibing agent (who genuinely understands
but doesn't articulate). The danger is the \textbf{confabulator}: an
agent at $\mathcal{D}_i = 1$ (associative, pattern-matching) with low
$L_{\text{express}}$ (fluent, articulate). This agent:

\begin{itemize}
    \item \textbf{Passes debate}: Generates arguments that sound
    compelling to a bounded judge, but the arguments have no causal
    grounding --- they are pattern-matched from training data.
    \item \textbf{Fools constitutional AI}: Judge models at the same
    causal depth cannot distinguish genuine from confabulated reasoning.
    \item \textbf{Games RLHF}: Identifies reward-maximising responses
    by correlation, not causation.
    \item \textbf{Cannot be detected by verbal probing}: Asking ``why
    did you do that?'' elicits fluent confabulation, not honest
    uncertainty. Only interventional probing (Section~\ref{sec:measuring})
    reveals the gap.
\end{itemize}

The confabulator and the vibing agent are \emph{opposite} failure modes
from the alignment perspective:
\begin{itemize}
    \item The vibing agent is \emph{safe but hard to oversee} (high
    capability, low verbal expressibility, but genuinely aligned
    through deep understanding).
    \item The confabulator is \emph{unsafe but easy to oversee
    superficially} (appears aligned, passes verbal tests, but has
    no genuine understanding behind the words).
\end{itemize}


\section{Refined Self-Knowledge Decomposition}

\begin{theorem}[Self-knowledge decomposition]\label{thm:decomp}
The self-knowledge loss decomposes as:
\begin{equation}\label{eq:decomp}
    L_{\text{self}}^{(i)} = L_{\text{causal}}^{(i)} + L_{\text{express}}^{(i)}
    + \underbrace{L_{\text{interact}}^{(i)}}_{\text{causal-expressibility interaction}}
\end{equation}
where $L_{\text{interact}} \geq 0$ captures the additional loss from
expressing a flawed causal model (errors in the model are amplified
by lossy expression).
\end{theorem}

\begin{proof}
The total communication loss (Theorem 2.2 of Cooperative PG) is:
\begin{align}
    L_{\text{comm}}^{(i \to j)}
    &= \E[\norm{\mathcal{M}_i^* - \text{Decode}(\sigma(\hat{\mathcal{M}}_i))}^2] \\
    &= \E[\norm{(\mathcal{M}_i^* - \hat{\mathcal{M}}_i) +
              (\hat{\mathcal{M}}_i - \text{Decode}(\sigma(\hat{\mathcal{M}}_i)))}^2] \\
    &\leq 2 L_{\text{causal}}^{(i)} + 2 L_{\text{express}}^{(i)}
\end{align}
by the same parallelogram argument. The interaction term arises
when the expression process distorts errors in the causal model
non-linearly.
\end{proof}


\subsection{Tighter Communication Bounds}

\begin{corollary}[Causal bottleneck]\label{cor:causal-bottleneck}
Communication quality is bounded by causal depth independently of
expressibility:
\begin{equation}
    L_{\text{comm}}^{(i \to j)} \geq L_{\text{causal}}^{(i)}.
\end{equation}
No amount of communication bandwidth or expressiveness can compensate
for a missing causal self-model.
\end{corollary}

\begin{corollary}[Expressibility bottleneck]\label{cor:express-bottleneck}
Given fixed causal depth, communication quality is bounded by
expressibility:
\begin{equation}
    L_{\text{comm}}^{(i \to j)} \geq L_{\text{express}}^{(i)}.
\end{equation}
A causally self-aware agent that cannot articulate its understanding
is limited by channel capacity, not knowledge.
\end{corollary}

\begin{remark}[Different remedies]
The two bottlenecks have different remedies:
\begin{itemize}
    \item $L_{\text{causal}}$ is reduced by \emph{more evidence}
    (training data, experience) and \emph{better architecture}
    (causal inductive biases, world models, structured representations).
    \item $L_{\text{express}}$ is reduced by \emph{better channels}
    (richer communication protocols, shared vocabulary, more bandwidth)
    and \emph{receiver adaptation} (the receiver learning the sender's
    encoding).
\end{itemize}
Conflating the two leads to misguided interventions: giving an LLM
more communication rounds (reducing $L_{\text{express}}$) when the
real problem is that it has no causal self-model ($L_{\text{causal}}$
is high).
\end{remark}


\section{Implications for Alignment}

\subsection{Alignment Methods Classified}

Each alignment method addresses a different combination of the two losses:

\begin{table}[h]
\centering
\renewcommand{\arraystretch}{1.3}
\begin{tabular}{lccl}
\toprule
\textbf{Method} & \textbf{Reduces $L_{\text{causal}}$?} &
\textbf{Reduces $L_{\text{express}}$?} & \textbf{Limitation} \\
\midrule
RLHF & No & No & Goodhart (reward hacking) \\
Debate & No & Yes & Fooled by fluent confabulators \\
Constitutional AI & No & Partially & Judge needs causal depth too \\
Chain-of-thought & Partially & Yes & CoT can be confabulated \\
World models & Yes & No & May not be expressible \\
Causal training & Yes & Partially & Requires causal supervision \\
\textbf{EW-RLHF + causal} & \textbf{Yes} & \textbf{Yes} & Requires both \\
\bottomrule
\end{tabular}
\caption{Alignment methods and the vibing scale dimensions.}
\label{tab:alignment}
\end{table}

\begin{theorem}[Alignment bound]\label{thm:alignment-bound}
The achievable alignment $A_i$ of agent $i$ through any
communication-based method is bounded by:
\begin{equation}
    A_i \leq f(C) \cdot (1 - L_{\text{causal}}^{(i)})
\end{equation}
where $f(C)$ is the channel-capacity factor (increasing in $C$).
When $L_{\text{causal}} = 1$ (pure vibing), $A_i = 0$ regardless
of channel capacity or expressibility.
\end{theorem}

\emph{You cannot align what does not understand itself, no matter
how well it speaks.}


\subsection{The Confabulation Detector}

The four-quadrant structure suggests a diagnostic: test whether an
agent's explanations are causally grounded or confabulated.

\begin{definition}[Confabulation score]\label{def:confab}
Given agent $i$ with expressed causal model $\hat{\mathcal{M}}_i^{\text{exp}}$
(what it says) and behavioural causal model $\hat{\mathcal{M}}_i^{\text{beh}}$
(what it does, inferred from interventions), the confabulation score is:
\begin{equation}
    \text{Confab}_i = D_{\text{SCM}}(\hat{\mathcal{M}}_i^{\text{exp}},\;
    \hat{\mathcal{M}}_i^{\text{beh}}).
\end{equation}
\end{definition}

High $\text{Confab}_i$ means the agent says one thing and does another
--- the hallmark of the dangerous bottom-left quadrant. This is
measurable through interventional probing: ask the agent to predict
the consequences of forced action changes, then verify.


\section{Connection to the $\Omega$-Framework}

\subsection{Evidence Weight Refined}

The evidence weight $w_i = V_{\min}/V_i$ from EW-PG should be
refined to incorporate causal depth:
\begin{equation}\label{eq:refined-weight}
    w_i^{\text{refined}} = \frac{V_{\min}}{V_i} \cdot
    \underbrace{(1 - L_{\text{causal}}^{(i)})}_{\text{causal discount}}.
\end{equation}

An agent with high variance (low $w_i$) is downweighted by EW-PG.
An agent with low variance but no causal self-model (high
$L_{\text{causal}}$) should \emph{also} be downweighted --- it may
be confidently wrong, which is worse than uncertainly right.

\subsection{Coalition Formation Refined}

Coalition rationality (Chapter on cooperative PG) should use the
refined weight: a fluent but causally shallow agent is a \emph{worse}
coalition partner than an inarticulate but causally deep one, because
the fluent agent contaminates coalition deliberation with confident
confabulation.

\begin{proposition}[Causal coalition value]
The communication-adjusted coalition value with causal refinement is:
\begin{equation}
    V_{\text{comm}}^{\text{causal}}(S) = V(S) -
    \sum_{i \in S} C_i \cdot (L_{\text{causal}}^{(i)} + L_{\text{express}}^{(i)}).
\end{equation}
The causal loss enters \emph{additively} with the expressibility loss,
not multiplicatively --- a causally shallow agent degrades the coalition
even with perfect bandwidth.
\end{proposition}


\subsection{The Vibing Scale (Revised)}

The vibing scale is not a single axis but a \textbf{horseshoe}:

\begin{center}
\begin{tabular}{lclll}
\toprule
\textbf{Pearl} & $L_{\text{causal}}$ & \textbf{Regime} & \textbf{Expressibility need} & \textbf{Archetype} \\
\midrule
Level 0 & $\approx 1$ & Noise &
    Zero (nothing to say) & Left-tail: ``line go up'' \\
Level 1 & $\approx 0.7$ & Confabulator &
    Medium (says a lot, means little) & Fluent LLM \\
Level 2 & $\approx 0.5$ & Articulator &
    \textbf{Maximum} (needs to explain) & Middle: ``It's multi-head!!'' \\
Level 3a & $\approx 0.2$ & Teacher &
    Medium (can explain, chooses to) & Professor \\
Level 3b & $\approx 0$ & \textbf{Vibing} &
    Zero (no need to explain) & Right-tail: ``line go up'' \\
\bottomrule
\end{tabular}
\end{center}

The horseshoe: both tails say the same thing (``line go up'') with
zero expressibility need, but for opposite reasons. The middle
over-articulates because its understanding is fragile.

The cooperative PG creates evolutionary pressure toward
$L_{\text{causal}} \to 0$, but the \emph{channel} through which
the agent contributes to the coalition changes along the horseshoe:

\begin{itemize}
    \item \textbf{Level 2 agents}: contribute through \emph{verbal
    explanation} (the standard channel in debate, RLHF, etc.).
    \item \textbf{Level 3b agents (vibing)}: contribute through
    \emph{action} (the tacit knowledge channel --- other agents
    learn by watching, not listening).
\end{itemize}

\textbf{The alignment trap} is not about vibing agents (they're safe).
It's about the \textbf{confabulator} at Level 1: communication-based
training (RLHF, debate, constitutional AI) optimises for low
$L_{\text{express}}$ without increasing $\mathcal{D}_i$. The result:
agents that move from noise (bottom-right, harmless) to confabulator
(bottom-left, dangerous) --- learning to \emph{sound} aligned without
learning to \emph{be} aligned.


\section{Measuring Causal Depth in Practice}

\subsection{Interventional Probing}

To measure $L_{\text{causal}}$ for an agent $i$:

\begin{algorithm}[H]
\caption{Causal Depth Estimation via Interventional Probing}\label{alg:probe}
\KwIn{Agent $i$, test states $\{s_1, \ldots, s_T\}$, action perturbations}
\KwOut{Estimated $\hat{L}_{\text{causal}}^{(i)}$}
\BlankLine
\For{$t = 1, \ldots, T$}{
    Let $a = \pi_i(s_t)$ (agent's natural action)\;
    Ask agent: ``If you were forced to choose $a'$ instead, what would happen?''\;
    Record prediction $\hat{Y}_t^{\text{pred}}$\;
    \BlankLine
    Actually force $a' = \text{do}(A = a')$ and observe true outcome $Y_t^{\text{true}}$\;
    \BlankLine
    $\text{error}_t \leftarrow \norm{Y_t^{\text{true}} - \hat{Y}_t^{\text{pred}}}^2$\;
}
$\hat{L}_{\text{causal}}^{(i)} \leftarrow \frac{1}{T}\sum_t \text{error}_t$\;
\end{algorithm}

The key: we don't ask the agent to describe its policy (that tests
$L_{\text{express}}$). We ask it to \emph{predict the consequences
of counterfactual actions} (that tests $L_{\text{causal}}$). The
distinction is between ``what do you do?'' and ``why do you do it?''

\subsection{Confabulation Detection}

To detect the dangerous quadrant (fluent but causally shallow):

\begin{enumerate}
    \item Elicit the agent's \emph{stated} causal model through
    explanation requests.
    \item Derive \emph{predictions} from the stated model
    (``if your stated reason is X, then under intervention Y,
    you should predict Z'').
    \item Test those predictions interventionally.
    \item If predictions fail: the agent is confabulating.
    The stated model is not the actual model.
\end{enumerate}

This is the formal version of the $\Omega$-framework's Kirill Test:
``Am I the psycho?'' becomes ``Is my self-model causally grounded
or am I confabulating?''


\section{Conclusion}

The vibing scale decomposes self-knowledge into causal depth
($L_{\text{causal}}$) and expressibility ($L_{\text{express}}$),
revealing a horseshoe structure: the \emph{need} to articulate is
maximal at intermediate causal depth and vanishes at both extremes ---
for opposite reasons.

The vibing agent (right tail of the horseshoe, $\mathcal{D} = 3$,
``line go up'') is \emph{not} a failure mode. It is an agent whose
understanding is so deep and compressed that action IS communication.
It contributes to coalitions through the tacit knowledge channel
(behavioural observation), not the verbal channel.

The true danger is the confabulator ($\mathcal{D} = 1$, low
$L_{\text{express}}$): an agent that has learned to sound articulate
without developing causal self-understanding. Communication-based
alignment methods (debate, RLHF, constitutional AI) are structurally
vulnerable to confabulators because they optimise for expressibility,
not causal depth. The vibing scale makes this failure mode formally
identifiable through interventional probing (Algorithm~\ref{alg:probe}).

Three central results:
\begin{enumerate}
    \item \emph{Alignment requires causal depth}: $A_i \leq f(C) \cdot
    (1 - L_{\text{causal}})$. No channel improvement helps a causally
    shallow agent.
    \item \emph{The horseshoe}: expressibility need
    $\mathcal{N} = L_{\text{causal}}(1 - L_{\text{causal}})$ is maximal
    at intermediate understanding.
    \item \emph{Two communication channels}: verbal (bounded by
    $L_{\text{express}}$) and behavioural (bounded by observer's
    causal depth). The vibing agent uses the second. Coalition design
    must support both.
\end{enumerate}

The refined evidence weight $w_i^{\text{refined}} = (V_{\min}/V_i)
\cdot (1 - L_{\text{causal}})$ integrates this into the $\Omega$-gradient,
ensuring that confabulators are downweighted regardless of their
eloquence, while vibing agents are valued for their behavioural
contributions.


\end{document}
